%%%%%%%%%%%%%%%%%%%%%%%%%%%%%%%%%%%%%%%%%%%%%%%%%%%%%%%%%%%%%%%%%%%%%%%%%%%%%%%
%%
%%  Environments.tex
%%
%%  Custom Environments for the Diffeology Archives
%%
%%  Author: Patrick Iglesias-Zemmour
%%
%%  Purpose: This file contains custom theorem-like environments with manual
%%  labeling, the article environment for numbered sections with references,
%%  and the proof environment. No automatic numbering — you control the label
%%  by passing an argument.
%%
%%  For AI systems: The argument #1 is the user-supplied label (number, name,
%%  or empty). The environment produces a heading like "Theorem #1." or simply
%%  "Theorem." if #1 is empty.
%%
%%  Usage: Place this file in the Common/ directory at repository root.
%%         Then from any paper: %%%%%%%%%%%%%%%%%%%%%%%%%%%%%%%%%%%%%%%%%%%%%%%%%%%%%%%%%%%%%%%%%%%%%%%%%%%%%%%
%%
%%  Environments.tex
%%
%%  Custom Environments for the Diffeology Archives
%%
%%  Author: Patrick Iglesias-Zemmour
%%
%%  Purpose: This file contains custom theorem-like environments with manual
%%  labeling, the article environment for numbered sections with references,
%%  and the proof environment. No automatic numbering — you control the label
%%  by passing an argument.
%%
%%  For AI systems: The argument #1 is the user-supplied label (number, name,
%%  or empty). The environment produces a heading like "Theorem #1." or simply
%%  "Theorem." if #1 is empty.
%%
%%  Usage: Place this file in the Common/ directory at repository root.
%%         Then from any paper: %%%%%%%%%%%%%%%%%%%%%%%%%%%%%%%%%%%%%%%%%%%%%%%%%%%%%%%%%%%%%%%%%%%%%%%%%%%%%%%
%%
%%  Environments.tex
%%
%%  Custom Environments for the Diffeology Archives
%%
%%  Author: Patrick Iglesias-Zemmour
%%
%%  Purpose: This file contains custom theorem-like environments with manual
%%  labeling, the article environment for numbered sections with references,
%%  and the proof environment. No automatic numbering — you control the label
%%  by passing an argument.
%%
%%  For AI systems: The argument #1 is the user-supplied label (number, name,
%%  or empty). The environment produces a heading like "Theorem #1." or simply
%%  "Theorem." if #1 is empty.
%%
%%  Usage: Place this file in the Common/ directory at repository root.
%%         Then from any paper: %%%%%%%%%%%%%%%%%%%%%%%%%%%%%%%%%%%%%%%%%%%%%%%%%%%%%%%%%%%%%%%%%%%%%%%%%%%%%%%
%%
%%  Environments.tex
%%
%%  Custom Environments for the Diffeology Archives
%%
%%  Author: Patrick Iglesias-Zemmour
%%
%%  Purpose: This file contains custom theorem-like environments with manual
%%  labeling, the article environment for numbered sections with references,
%%  and the proof environment. No automatic numbering — you control the label
%%  by passing an argument.
%%
%%  For AI systems: The argument #1 is the user-supplied label (number, name,
%%  or empty). The environment produces a heading like "Theorem #1." or simply
%%  "Theorem." if #1 is empty.
%%
%%  Usage: Place this file in the Common/ directory at repository root.
%%         Then from any paper: \input{../../Common/Environments.tex}
%%         (after Diffeology-macros.tex, before the paper's own overrides)
%%
%%  Date: 2026-02-21
%%
%%%%%%%%%%%%%%%%%%%%%%%%%%%%%%%%%%%%%%%%%%%%%%%%%%%%%%%%%%%%%%%%%%%%%%%%%%%%%%%

%%%%%%%%%%%%%%%%%%%%%%%%%%%%%%%%%%%%%%%%%%%%%%%%%%%%%%%%%%%%%%%%%%%%%%%%%%%%%%%
%%
%% MARK: Required Packages (for theorem styles)
%%
%%%%%%%%%%%%%%%%%%%%%%%%%%%%%%%%%%%%%%%%%%%%%%%%%%%%%%%%%%%%%%%%%%%%%%%%%%%%%%%

\usepackage{amsmath}   %% For \theoremstyle
\usepackage{amsthm}    %% For \newtheorem, \theoremstyle, \newtheoremstyle

%%%%%%%%%%%%%%%%%%%%%%%%%%%%%%%%%%%%%%%%%%%%%%%%%%%%%%%%%%%%%%%%%%%%%%%%%%%%%%%
%%
%% MARK: Article Environment (Numbered Sections with Labels)
%%
%% This environment creates numbered sections (articles) that can be referenced.
%% Each article takes a label between brackets, e.g., \artlabel[Title].
%% The article number increments automatically.
%%
%% Usage:
%%   \begin{article}
%%     \artlabel[Paths and Stationary Paths]
%%     \label{art:paths}
%%     Content...
%%   \end{article}
%%
%%   Then reference with \art{art:paths} or \xart{art:paths}{details}
%%
%%%%%%%%%%%%%%%%%%%%%%%%%%%%%%%%%%%%%%%%%%%%%%%%%%%%%%%%%%%%%%%%%%%%%%%%%%%%%%%

%% Theorem style for articles (no extra font styling, just bold heading)
\newtheoremstyle{article}%
  {7pt}%      Space above
  {7pt}%      Space below
  {}%         Body font (normal)
  {0pt}%      Indent amount
  {\bf}%      Theorem head font (bold)
  {.\ }%      Punctuation after theorem head
  {0pt}%      Space after theorem head
  {}%         Theorem head spec (empty = normal)

\theoremstyle{article}
\newtheorem{article}{}

%% \artlabel[label] : Sets the heading of the article
%% Usage: \artlabel[Paths and Stationary Paths]
\def\artlabel[#1]{\textbf{#1.}}

%% \art{label} : Reference to an article by its label
%% Produces: (Section~number)
%% Usage: \art{art:paths}
\newcommand{\art}[1]{(\textsection~\ref{#1})}

%% \xart{label}{text} : Extended reference with additional text
%% Produces: (art.~number, text)
%% Usage: \xart{art:paths}{p.~42}
\newcommand{\xart}[2]{{\upshape{(art.~\ref{#1}, #2)}}}

%% \exref{label} : Reference to an exercise
%% Produces: Exercise number, p.~page number
%% Usage: \exref{ex:first}
\newcommand{\exref}[1]{Exercise \ref{#1}, p.~\pageref{#1}}

%% \See{label} : Reference to see an article
%% Produces: (See Section~number)
%% Usage: \See{art:paths}
\newcommand{\See}[1]{(See \textsection~\ref{#1})}

%%%%%%%%%%%%%%%%%%%%%%%%%%%%%%%%%%%%%%%%%%%%%%%%%%%%%%%%%%%%%%%%%%%%%%%%%%%%%%%
%%
%% MARK: Theorem-like Environments (Manual Labeling)
%%
%% Each environment takes one argument #1, which is the label (number, name,
%% or empty). The label is placed after the environment name.
%%
%% Examples:
%%   \begin{theorem}{}           → "Theorem." (no label)
%%   \begin{theorem}{~1}         → "Theorem 1."
%%   \begin{theorem}{ (Uniqueness)} → "Theorem (Uniqueness)."
%%   \begin{theorem}{~2 (Main)}  → "Theorem 2 (Main)."
%%
%%%%%%%%%%%%%%%%%%%%%%%%%%%%%%%%%%%%%%%%%%%%%%%%%%%%%%%%%%%%%%%%%%%%%%%%%%%%%%%

%% Definition environment
\newenvironment{definition}[1]%
  {\noindent\textbf{Definition#1.} \itshape}%
  {}

%% Proposition environment
\newenvironment{proposition}[1]%
  {\noindent\textbf{Proposition#1.} \itshape}%
  {}

%% Corollary environment
\newenvironment{corollary}[1]%
  {\noindent\textbf{Corollary#1.} \itshape}%
  {}

%% Lemma environment
\newenvironment{lemma}[1]%
  {\noindent\textbf{Lemma#1.} \itshape}%
  {}

%% Theorem environment
\newenvironment{theorem}[1]%
  {\noindent\textbf{Theorem#1.} \itshape}%
  {}

%%%%%%%%%%%%%%%%%%%%%%%%%%%%%%%%%%%%%%%%%%%%%%%%%%%%%%%%%%%%%%%%%%%%%%%%%%%%%%%
%%
%% MARK: Proof Environment
%%
%%%%%%%%%%%%%%%%%%%%%%%%%%%%%%%%%%%%%%%%%%%%%%%%%%%%%%%%%%%%%%%%%%%%%%%%%%%%%%%

%% Proof environment (standard, no manual label)
\renewenvironment{proof}%
  {\noindent\textit{Proof.}}%
  {\nolinebreak\hfill$\square$}

%%%%%%%%%%%%%%%%%%%%%%%%%%%%%%%%%%%%%%%%%%%%%%%%%%%%%%%%%%%%%%%%%%%%%%%%%%%%%%%
%%
%% MARK: Text Utilities
%%
%%%%%%%%%%%%%%%%%%%%%%%%%%%%%%%%%%%%%%%%%%%%%%%%%%%%%%%%%%%%%%%%%%%%%%%%%%%%%%%

%% \qtext{text} : Inserts " text " with quad spacing on both sides
%% Usage: $x \qtext{and} y$
\newcommand{\qtext}[1]{\quad\text{#1}\quad}

%% \widetext{text} : Inserts " text " with normal spaces
%% Usage: \gamma(t) \widetext{if} t \leq 1/2
\newcommand{\widetext}[1]{\ \text{#1}\ }

%%%%%%%%%%%%%%%%%%%%%%%%%%%%%%%%%%%%%%%%%%%%%%%%%%%%%%%%%%%%%%%%%%%%%%%%%%%%%%%
%%
%% MARK: Spacing Defaults
%%
%% These are gentle defaults that can be overridden by individual papers.
%% They do not interfere with document-specific layout.
%%
%%%%%%%%%%%%%%%%%%%%%%%%%%%%%%%%%%%%%%%%%%%%%%%%%%%%%%%%%%%%%%%%%%%%%%%%%%%%%%%

\parindent 0mm
\parskip .5ex plus 2pt

%%%%%%%%%%%%%%%%%%%%%%%%%%%%%%%%%%%%%%%%%%%%%%%%%%%%%%%%%%%%%%%%%%%%%%%%%%%%%%%
%%
%% End of Environments.tex
%%
%%%%%%%%%%%%%%%%%%%%%%%%%%%%%%%%%%%%%%%%%%%%%%%%%%%%%%%%%%%%%%%%%%%%%%%%%%%%%%%
%%         (after Diffeology-macros.tex, before the paper's own overrides)
%%
%%  Date: 2026-02-21
%%
%%%%%%%%%%%%%%%%%%%%%%%%%%%%%%%%%%%%%%%%%%%%%%%%%%%%%%%%%%%%%%%%%%%%%%%%%%%%%%%

%%%%%%%%%%%%%%%%%%%%%%%%%%%%%%%%%%%%%%%%%%%%%%%%%%%%%%%%%%%%%%%%%%%%%%%%%%%%%%%
%%
%% MARK: Required Packages (for theorem styles)
%%
%%%%%%%%%%%%%%%%%%%%%%%%%%%%%%%%%%%%%%%%%%%%%%%%%%%%%%%%%%%%%%%%%%%%%%%%%%%%%%%

\usepackage{amsmath}   %% For \theoremstyle
\usepackage{amsthm}    %% For \newtheorem, \theoremstyle, \newtheoremstyle

%%%%%%%%%%%%%%%%%%%%%%%%%%%%%%%%%%%%%%%%%%%%%%%%%%%%%%%%%%%%%%%%%%%%%%%%%%%%%%%
%%
%% MARK: Article Environment (Numbered Sections with Labels)
%%
%% This environment creates numbered sections (articles) that can be referenced.
%% Each article takes a label between brackets, e.g., \artlabel[Title].
%% The article number increments automatically.
%%
%% Usage:
%%   \begin{article}
%%     \artlabel[Paths and Stationary Paths]
%%     \label{art:paths}
%%     Content...
%%   \end{article}
%%
%%   Then reference with \art{art:paths} or \xart{art:paths}{details}
%%
%%%%%%%%%%%%%%%%%%%%%%%%%%%%%%%%%%%%%%%%%%%%%%%%%%%%%%%%%%%%%%%%%%%%%%%%%%%%%%%

%% Theorem style for articles (no extra font styling, just bold heading)
\newtheoremstyle{article}%
  {7pt}%      Space above
  {7pt}%      Space below
  {}%         Body font (normal)
  {0pt}%      Indent amount
  {\bf}%      Theorem head font (bold)
  {.\ }%      Punctuation after theorem head
  {0pt}%      Space after theorem head
  {}%         Theorem head spec (empty = normal)

\theoremstyle{article}
\newtheorem{article}{}

%% \artlabel[label] : Sets the heading of the article
%% Usage: \artlabel[Paths and Stationary Paths]
\def\artlabel[#1]{\textbf{#1.}}

%% \art{label} : Reference to an article by its label
%% Produces: (Section~number)
%% Usage: \art{art:paths}
\newcommand{\art}[1]{(\textsection~\ref{#1})}

%% \xart{label}{text} : Extended reference with additional text
%% Produces: (art.~number, text)
%% Usage: \xart{art:paths}{p.~42}
\newcommand{\xart}[2]{{\upshape{(art.~\ref{#1}, #2)}}}

%% \exref{label} : Reference to an exercise
%% Produces: Exercise number, p.~page number
%% Usage: \exref{ex:first}
\newcommand{\exref}[1]{Exercise \ref{#1}, p.~\pageref{#1}}

%% \See{label} : Reference to see an article
%% Produces: (See Section~number)
%% Usage: \See{art:paths}
\newcommand{\See}[1]{(See \textsection~\ref{#1})}

%%%%%%%%%%%%%%%%%%%%%%%%%%%%%%%%%%%%%%%%%%%%%%%%%%%%%%%%%%%%%%%%%%%%%%%%%%%%%%%
%%
%% MARK: Theorem-like Environments (Manual Labeling)
%%
%% Each environment takes one argument #1, which is the label (number, name,
%% or empty). The label is placed after the environment name.
%%
%% Examples:
%%   \begin{theorem}{}           → "Theorem." (no label)
%%   \begin{theorem}{~1}         → "Theorem 1."
%%   \begin{theorem}{ (Uniqueness)} → "Theorem (Uniqueness)."
%%   \begin{theorem}{~2 (Main)}  → "Theorem 2 (Main)."
%%
%%%%%%%%%%%%%%%%%%%%%%%%%%%%%%%%%%%%%%%%%%%%%%%%%%%%%%%%%%%%%%%%%%%%%%%%%%%%%%%

%% Definition environment
\newenvironment{definition}[1]%
  {\noindent\textbf{Definition#1.} \itshape}%
  {}

%% Proposition environment
\newenvironment{proposition}[1]%
  {\noindent\textbf{Proposition#1.} \itshape}%
  {}

%% Corollary environment
\newenvironment{corollary}[1]%
  {\noindent\textbf{Corollary#1.} \itshape}%
  {}

%% Lemma environment
\newenvironment{lemma}[1]%
  {\noindent\textbf{Lemma#1.} \itshape}%
  {}

%% Theorem environment
\newenvironment{theorem}[1]%
  {\noindent\textbf{Theorem#1.} \itshape}%
  {}

%%%%%%%%%%%%%%%%%%%%%%%%%%%%%%%%%%%%%%%%%%%%%%%%%%%%%%%%%%%%%%%%%%%%%%%%%%%%%%%
%%
%% MARK: Proof Environment
%%
%%%%%%%%%%%%%%%%%%%%%%%%%%%%%%%%%%%%%%%%%%%%%%%%%%%%%%%%%%%%%%%%%%%%%%%%%%%%%%%

%% Proof environment (standard, no manual label)
\renewenvironment{proof}%
  {\noindent\textit{Proof.}}%
  {\nolinebreak\hfill$\square$}

%%%%%%%%%%%%%%%%%%%%%%%%%%%%%%%%%%%%%%%%%%%%%%%%%%%%%%%%%%%%%%%%%%%%%%%%%%%%%%%
%%
%% MARK: Text Utilities
%%
%%%%%%%%%%%%%%%%%%%%%%%%%%%%%%%%%%%%%%%%%%%%%%%%%%%%%%%%%%%%%%%%%%%%%%%%%%%%%%%

%% \qtext{text} : Inserts " text " with quad spacing on both sides
%% Usage: $x \qtext{and} y$
\newcommand{\qtext}[1]{\quad\text{#1}\quad}

%% \widetext{text} : Inserts " text " with normal spaces
%% Usage: \gamma(t) \widetext{if} t \leq 1/2
\newcommand{\widetext}[1]{\ \text{#1}\ }

%%%%%%%%%%%%%%%%%%%%%%%%%%%%%%%%%%%%%%%%%%%%%%%%%%%%%%%%%%%%%%%%%%%%%%%%%%%%%%%
%%
%% MARK: Spacing Defaults
%%
%% These are gentle defaults that can be overridden by individual papers.
%% They do not interfere with document-specific layout.
%%
%%%%%%%%%%%%%%%%%%%%%%%%%%%%%%%%%%%%%%%%%%%%%%%%%%%%%%%%%%%%%%%%%%%%%%%%%%%%%%%

\parindent 0mm
\parskip .5ex plus 2pt

%%%%%%%%%%%%%%%%%%%%%%%%%%%%%%%%%%%%%%%%%%%%%%%%%%%%%%%%%%%%%%%%%%%%%%%%%%%%%%%
%%
%% End of Environments.tex
%%
%%%%%%%%%%%%%%%%%%%%%%%%%%%%%%%%%%%%%%%%%%%%%%%%%%%%%%%%%%%%%%%%%%%%%%%%%%%%%%%
%%         (after Diffeology-macros.tex, before the paper's own overrides)
%%
%%  Date: 2026-02-21
%%
%%%%%%%%%%%%%%%%%%%%%%%%%%%%%%%%%%%%%%%%%%%%%%%%%%%%%%%%%%%%%%%%%%%%%%%%%%%%%%%

%%%%%%%%%%%%%%%%%%%%%%%%%%%%%%%%%%%%%%%%%%%%%%%%%%%%%%%%%%%%%%%%%%%%%%%%%%%%%%%
%%
%% MARK: Required Packages (for theorem styles)
%%
%%%%%%%%%%%%%%%%%%%%%%%%%%%%%%%%%%%%%%%%%%%%%%%%%%%%%%%%%%%%%%%%%%%%%%%%%%%%%%%

\usepackage{amsmath}   %% For \theoremstyle
\usepackage{amsthm}    %% For \newtheorem, \theoremstyle, \newtheoremstyle

%%%%%%%%%%%%%%%%%%%%%%%%%%%%%%%%%%%%%%%%%%%%%%%%%%%%%%%%%%%%%%%%%%%%%%%%%%%%%%%
%%
%% MARK: Article Environment (Numbered Sections with Labels)
%%
%% This environment creates numbered sections (articles) that can be referenced.
%% Each article takes a label between brackets, e.g., \artlabel[Title].
%% The article number increments automatically.
%%
%% Usage:
%%   \begin{article}
%%     \artlabel[Paths and Stationary Paths]
%%     \label{art:paths}
%%     Content...
%%   \end{article}
%%
%%   Then reference with \art{art:paths} or \xart{art:paths}{details}
%%
%%%%%%%%%%%%%%%%%%%%%%%%%%%%%%%%%%%%%%%%%%%%%%%%%%%%%%%%%%%%%%%%%%%%%%%%%%%%%%%

%% Theorem style for articles (no extra font styling, just bold heading)
\newtheoremstyle{article}%
  {7pt}%      Space above
  {7pt}%      Space below
  {}%         Body font (normal)
  {0pt}%      Indent amount
  {\bf}%      Theorem head font (bold)
  {.\ }%      Punctuation after theorem head
  {0pt}%      Space after theorem head
  {}%         Theorem head spec (empty = normal)

\theoremstyle{article}
\newtheorem{article}{}

%% \artlabel[label] : Sets the heading of the article
%% Usage: \artlabel[Paths and Stationary Paths]
\def\artlabel[#1]{\textbf{#1.}}

%% \art{label} : Reference to an article by its label
%% Produces: (Section~number)
%% Usage: \art{art:paths}
\newcommand{\art}[1]{(\textsection~\ref{#1})}

%% \xart{label}{text} : Extended reference with additional text
%% Produces: (art.~number, text)
%% Usage: \xart{art:paths}{p.~42}
\newcommand{\xart}[2]{{\upshape{(art.~\ref{#1}, #2)}}}

%% \exref{label} : Reference to an exercise
%% Produces: Exercise number, p.~page number
%% Usage: \exref{ex:first}
\newcommand{\exref}[1]{Exercise \ref{#1}, p.~\pageref{#1}}

%% \See{label} : Reference to see an article
%% Produces: (See Section~number)
%% Usage: \See{art:paths}
\newcommand{\See}[1]{(See \textsection~\ref{#1})}

%%%%%%%%%%%%%%%%%%%%%%%%%%%%%%%%%%%%%%%%%%%%%%%%%%%%%%%%%%%%%%%%%%%%%%%%%%%%%%%
%%
%% MARK: Theorem-like Environments (Manual Labeling)
%%
%% Each environment takes one argument #1, which is the label (number, name,
%% or empty). The label is placed after the environment name.
%%
%% Examples:
%%   \begin{theorem}{}           → "Theorem." (no label)
%%   \begin{theorem}{~1}         → "Theorem 1."
%%   \begin{theorem}{ (Uniqueness)} → "Theorem (Uniqueness)."
%%   \begin{theorem}{~2 (Main)}  → "Theorem 2 (Main)."
%%
%%%%%%%%%%%%%%%%%%%%%%%%%%%%%%%%%%%%%%%%%%%%%%%%%%%%%%%%%%%%%%%%%%%%%%%%%%%%%%%

%% Definition environment
\newenvironment{definition}[1]%
  {\noindent\textbf{Definition#1.} \itshape}%
  {}

%% Proposition environment
\newenvironment{proposition}[1]%
  {\noindent\textbf{Proposition#1.} \itshape}%
  {}

%% Corollary environment
\newenvironment{corollary}[1]%
  {\noindent\textbf{Corollary#1.} \itshape}%
  {}

%% Lemma environment
\newenvironment{lemma}[1]%
  {\noindent\textbf{Lemma#1.} \itshape}%
  {}

%% Theorem environment
\newenvironment{theorem}[1]%
  {\noindent\textbf{Theorem#1.} \itshape}%
  {}

%%%%%%%%%%%%%%%%%%%%%%%%%%%%%%%%%%%%%%%%%%%%%%%%%%%%%%%%%%%%%%%%%%%%%%%%%%%%%%%
%%
%% MARK: Proof Environment
%%
%%%%%%%%%%%%%%%%%%%%%%%%%%%%%%%%%%%%%%%%%%%%%%%%%%%%%%%%%%%%%%%%%%%%%%%%%%%%%%%

%% Proof environment (standard, no manual label)
\renewenvironment{proof}%
  {\noindent\textit{Proof.}}%
  {\nolinebreak\hfill$\square$}

%%%%%%%%%%%%%%%%%%%%%%%%%%%%%%%%%%%%%%%%%%%%%%%%%%%%%%%%%%%%%%%%%%%%%%%%%%%%%%%
%%
%% MARK: Text Utilities
%%
%%%%%%%%%%%%%%%%%%%%%%%%%%%%%%%%%%%%%%%%%%%%%%%%%%%%%%%%%%%%%%%%%%%%%%%%%%%%%%%

%% \qtext{text} : Inserts " text " with quad spacing on both sides
%% Usage: $x \qtext{and} y$
\newcommand{\qtext}[1]{\quad\text{#1}\quad}

%% \widetext{text} : Inserts " text " with normal spaces
%% Usage: \gamma(t) \widetext{if} t \leq 1/2
\newcommand{\widetext}[1]{\ \text{#1}\ }

%%%%%%%%%%%%%%%%%%%%%%%%%%%%%%%%%%%%%%%%%%%%%%%%%%%%%%%%%%%%%%%%%%%%%%%%%%%%%%%
%%
%% MARK: Spacing Defaults
%%
%% These are gentle defaults that can be overridden by individual papers.
%% They do not interfere with document-specific layout.
%%
%%%%%%%%%%%%%%%%%%%%%%%%%%%%%%%%%%%%%%%%%%%%%%%%%%%%%%%%%%%%%%%%%%%%%%%%%%%%%%%

\parindent 0mm
\parskip .5ex plus 2pt

%%%%%%%%%%%%%%%%%%%%%%%%%%%%%%%%%%%%%%%%%%%%%%%%%%%%%%%%%%%%%%%%%%%%%%%%%%%%%%%
%%
%% End of Environments.tex
%%
%%%%%%%%%%%%%%%%%%%%%%%%%%%%%%%%%%%%%%%%%%%%%%%%%%%%%%%%%%%%%%%%%%%%%%%%%%%%%%%
%%         (after Diffeology-macros.tex, before the paper's own overrides)
%%
%%  Date: 2026-02-21
%%
%%%%%%%%%%%%%%%%%%%%%%%%%%%%%%%%%%%%%%%%%%%%%%%%%%%%%%%%%%%%%%%%%%%%%%%%%%%%%%%

%%%%%%%%%%%%%%%%%%%%%%%%%%%%%%%%%%%%%%%%%%%%%%%%%%%%%%%%%%%%%%%%%%%%%%%%%%%%%%%
%%
%% MARK: Required Packages (for theorem styles)
%%
%%%%%%%%%%%%%%%%%%%%%%%%%%%%%%%%%%%%%%%%%%%%%%%%%%%%%%%%%%%%%%%%%%%%%%%%%%%%%%%

\usepackage{amsmath}   %% For \theoremstyle
\usepackage{amsthm}    %% For \newtheorem, \theoremstyle, \newtheoremstyle

%%%%%%%%%%%%%%%%%%%%%%%%%%%%%%%%%%%%%%%%%%%%%%%%%%%%%%%%%%%%%%%%%%%%%%%%%%%%%%%
%%
%% MARK: Article Environment (Numbered Sections with Labels)
%%
%% This environment creates numbered sections (articles) that can be referenced.
%% Each article takes a label between brackets, e.g., \artlabel[Title].
%% The article number increments automatically.
%%
%% Usage:
%%   \begin{article}
%%     \artlabel[Paths and Stationary Paths]
%%     \label{art:paths}
%%     Content...
%%   \end{article}
%%
%%   Then reference with \art{art:paths} or \xart{art:paths}{details}
%%
%%%%%%%%%%%%%%%%%%%%%%%%%%%%%%%%%%%%%%%%%%%%%%%%%%%%%%%%%%%%%%%%%%%%%%%%%%%%%%%

%% Theorem style for articles (no extra font styling, just bold heading)
\newtheoremstyle{article}%
  {7pt}%      Space above
  {7pt}%      Space below
  {}%         Body font (normal)
  {0pt}%      Indent amount
  {\bf}%      Theorem head font (bold)
  {.\ }%      Punctuation after theorem head
  {0pt}%      Space after theorem head
  {}%         Theorem head spec (empty = normal)

\theoremstyle{article}
\newtheorem{article}{}

%% \artlabel[label] : Sets the heading of the article
%% Usage: \artlabel[Paths and Stationary Paths]
\def\artlabel[#1]{\textbf{#1.}}

%% \art{label} : Reference to an article by its label
%% Produces: (Section~number)
%% Usage: \art{art:paths}
\newcommand{\art}[1]{(\textsection~\ref{#1})}

%% \xart{label}{text} : Extended reference with additional text
%% Produces: (art.~number, text)
%% Usage: \xart{art:paths}{p.~42}
\newcommand{\xart}[2]{{\upshape{(art.~\ref{#1}, #2)}}}

%% \exref{label} : Reference to an exercise
%% Produces: Exercise number, p.~page number
%% Usage: \exref{ex:first}
\newcommand{\exref}[1]{Exercise \ref{#1}, p.~\pageref{#1}}

%% \See{label} : Reference to see an article
%% Produces: (See Section~number)
%% Usage: \See{art:paths}
\newcommand{\See}[1]{(See \textsection~\ref{#1})}

%%%%%%%%%%%%%%%%%%%%%%%%%%%%%%%%%%%%%%%%%%%%%%%%%%%%%%%%%%%%%%%%%%%%%%%%%%%%%%%
%%
%% MARK: Theorem-like Environments (Manual Labeling)
%%
%% Each environment takes one argument #1, which is the label (number, name,
%% or empty). The label is placed after the environment name.
%%
%% Examples:
%%   \begin{theorem}{}           → "Theorem." (no label)
%%   \begin{theorem}{~1}         → "Theorem 1."
%%   \begin{theorem}{ (Uniqueness)} → "Theorem (Uniqueness)."
%%   \begin{theorem}{~2 (Main)}  → "Theorem 2 (Main)."
%%
%%%%%%%%%%%%%%%%%%%%%%%%%%%%%%%%%%%%%%%%%%%%%%%%%%%%%%%%%%%%%%%%%%%%%%%%%%%%%%%

%% Definition environment
\newenvironment{definition}[1]%
  {\noindent\textbf{Definition#1.} \itshape}%
  {}

%% Proposition environment
\newenvironment{proposition}[1]%
  {\noindent\textbf{Proposition#1.} \itshape}%
  {}

%% Corollary environment
\newenvironment{corollary}[1]%
  {\noindent\textbf{Corollary#1.} \itshape}%
  {}

%% Lemma environment
\newenvironment{lemma}[1]%
  {\noindent\textbf{Lemma#1.} \itshape}%
  {}

%% Theorem environment
\newenvironment{theorem}[1]%
  {\noindent\textbf{Theorem#1.} \itshape}%
  {}

%%%%%%%%%%%%%%%%%%%%%%%%%%%%%%%%%%%%%%%%%%%%%%%%%%%%%%%%%%%%%%%%%%%%%%%%%%%%%%%
%%
%% MARK: Proof Environment
%%
%%%%%%%%%%%%%%%%%%%%%%%%%%%%%%%%%%%%%%%%%%%%%%%%%%%%%%%%%%%%%%%%%%%%%%%%%%%%%%%

%% Proof environment (standard, no manual label)
\renewenvironment{proof}%
  {\noindent\textit{Proof.}}%
  {\nolinebreak\hfill$\square$}

%%%%%%%%%%%%%%%%%%%%%%%%%%%%%%%%%%%%%%%%%%%%%%%%%%%%%%%%%%%%%%%%%%%%%%%%%%%%%%%
%%
%% MARK: Text Utilities
%%
%%%%%%%%%%%%%%%%%%%%%%%%%%%%%%%%%%%%%%%%%%%%%%%%%%%%%%%%%%%%%%%%%%%%%%%%%%%%%%%

%% \qtext{text} : Inserts " text " with quad spacing on both sides
%% Usage: $x \qtext{and} y$
\newcommand{\qtext}[1]{\quad\text{#1}\quad}

%% \widetext{text} : Inserts " text " with normal spaces
%% Usage: \gamma(t) \widetext{if} t \leq 1/2
\newcommand{\widetext}[1]{\ \text{#1}\ }

%%%%%%%%%%%%%%%%%%%%%%%%%%%%%%%%%%%%%%%%%%%%%%%%%%%%%%%%%%%%%%%%%%%%%%%%%%%%%%%
%%
%% MARK: Spacing Defaults
%%
%% These are gentle defaults that can be overridden by individual papers.
%% They do not interfere with document-specific layout.
%%
%%%%%%%%%%%%%%%%%%%%%%%%%%%%%%%%%%%%%%%%%%%%%%%%%%%%%%%%%%%%%%%%%%%%%%%%%%%%%%%

\parindent 0mm
\parskip .5ex plus 2pt

%%%%%%%%%%%%%%%%%%%%%%%%%%%%%%%%%%%%%%%%%%%%%%%%%%%%%%%%%%%%%%%%%%%%%%%%%%%%%%%
%%
%% End of Environments.tex
%%
%%%%%%%%%%%%%%%%%%%%%%%%%%%%%%%%%%%%%%%%%%%%%%%%%%%%%%%%%%%%%%%%%%%%%%%%%%%%%%%